\documentclass[11pt]{article}

\usepackage[a4paper,margin=1in]{geometry}
\usepackage{amsmath,amssymb,amsthm}
\usepackage{microtype}
\usepackage[hidelinks]{hyperref}
\usepackage{listings}
\lstset{basicstyle=\ttfamily\small,columns=fullflexible,keepspaces=true,
  literate={ε}{{$\varepsilon$}}1 {≤}{{$\le$}}1 {ℕ}{{$\mathbb N$}}1 {ℝ}{{$\mathbb R$}}1}

\newtheorem{theorem}{Theorem}
\newtheorem{lemma}{Lemma}
\theoremstyle{remark}
\newtheorem*{remark}{Remark}
\DeclareMathOperator{\lcm}{lcm}

\title{The sharp constant in Erd\H{o}s Problem 457}
\author{Ji Ho Bae}
\date{October 2026}

\begin{document}
\maketitle

\begin{abstract}
Let \(q(n,k)\) be the least prime not dividing \((n+1)(n+2)\cdots(n+k)\).  We prove that
\(q(n,\log n)\ge(1-\varepsilon)\log n\,\log\log n/\log\log\log n\) for infinitely many \(n\),
improving the constant \(\tfrac12\) of Tao to the value \(1\) predicted by the probabilistic
model.  The proof is elementary and is formally verified in Lean~4.
\end{abstract}

\section{The result}

Erd\H{o}s and Pomerance \cite{Erdos1979} asked how large
\(q(n,\log n)\) (meaning \(q(n,\lfloor\log n\rfloor)\)) can be; see Problem~457 of
\cite{Bloom457}.  In the discussion of that problem, constructions found by GPT-5.2~Pro
(prompted by Barreto) showed \(q(n,\log n)\ge C\log n\) infinitely often for every \(C\), and
Tao showed
\[
q(n,\log n)>\bigl(\tfrac12-o(1)\bigr)\,g(n),\qquad
g(n):=\frac{\log n\,\log\log n}{\log\log\log n},
\]
for infinitely many \(n\), remarking that \(\tfrac12\) is the natural limit of his argument and
that \(g(n)\) is the expected order.

\begin{theorem}\label{thm}
For every \(\varepsilon>0\) there are infinitely many \(n\) with
\(q(n,\log n)\ge(1-\varepsilon)\,g(n)\).  Equivalently,
\(\limsup_{n\to\infty}q(n,\log n)/g(n)\ge1\).
\end{theorem}

The constant is best possible in the probabilistic model: for \(k=\lfloor\log X\rfloor\) and
\(Q=\lambda g(X)\), the expected number of \(n\in(X,2X]\) with \(q(n,k)>Q\) is
\(X\prod_{k<p\le Q}k/p=X^{1-\lambda+o(1)}\).  An upper bound of the form
\(q(n,\log n)<(1-c)(\log n)^2\) is Problem~1181 of \cite{Bloom457} and is open.

\section{Three lemmas}

Write \(\theta(x)=\sum_{p\le x}\log p\) and \(\psi(x)=\log\lcm(1,\dots,\lfloor x\rfloor)\).
Besides Chebyshev's bounds \(\theta(x)\le x\log4\) and \(\psi(x)\le C_1x\) with
\(C_1=\log4+4\), we use only the following.

\begin{lemma}[Chebyshev]\label{lem:cheb}
For every \(\delta>0\), \(\theta(x)\le(1+\delta)x\) for arbitrarily large \(x\).
\end{lemma}

\begin{proof}
Legendre's formula gives \(\sum_{p\le N}\lfloor N/p\rfloor\log p\le\log N!\le N\log N\), hence
\(S(N):=\sum_{p\le N}\frac{\log p}{p}\le\log N+\frac{\theta(N)}N\le\log N+\log4\).
Partial summation gives \(S(N)\ge\sum_{n<N}\theta(n)/(n(n+1))\).  If \(\theta(n)>(1+\delta)n\)
for all \(n\ge n_0\), the right side is \(\ge(1+\delta)\log N-O(1)\), a contradiction for large
\(N\).
\end{proof}

\begin{lemma}[Anchored pigeonhole]\label{lem:anchor}
Let \(\Phi:\{0,\dots,J\}\to\Omega\) with \(\Omega\) finite, and let \(F\) be a finite set of
positive integers with \(|\Omega|(|F|+1)<J+1\).  Then \(\Phi(i)=\Phi(j)\) for some
\(0\le i<j\le J\) with \(j-i\notin F\).
\end{lemma}

\begin{proof}
A fibre \(B\) of \(\Phi\) has more than \(|F|+1\) elements; with \(i=\min B\), the \(|B|-1>|F|\)
differences \(j-i\) \((j\in B\setminus\{i\})\) cannot all lie in \(F\).
\end{proof}

\begin{lemma}[Fold]\label{lem:fold}
Let \(2h<p\) and \(tL\equiv e\pmod p\) with \(|e|<h\).  Then \(p\) divides one of
\(n+2,\dots,n+2h\), where \(n=tL-h-1\).
\end{lemma}

\begin{proof}
\(n+(1+h-e)=tL-e\equiv0\pmod p\), and \(2\le1+h-e\le2h\).
\end{proof}

\section{Proof of Theorem~\ref{thm}}

Fix \(\varepsilon\in(0,\frac12]\) and put \(c=(3+C_1)/\varepsilon\).  For an integer \(M\) let
\[
A=\frac{c\log M}{\log\log M},\quad K=\lfloor AM\rfloor,\quad
h=\Bigl\lfloor\frac{M}{20\log A}\Bigr\rfloor,\quad L=\lcm(1,\dots,M),\quad
\mathcal P=\{p:M<p\le K\}.
\]

\emph{The code.}  Let \(\kappa(j)=(\lfloor[jL]_p/h\rfloor)_{p\in\mathcal P}\), where \([x]_p\)
is the least nonnegative residue.  It takes at most
\(C_M=\prod_{p\in\mathcal P}(\lfloor p/h\rfloor+1)\) values.  Since \(\log(p/h)/\log p\) increases
in \(p\),
\[
\log C_M\le\sum_{p\in\mathcal P}\Bigl(\log\frac ph+\frac hp\Bigr)
\le\theta(K)\,\frac{\log(K/h)}{\log K}+\pi(K)\frac hM .
\]
As \(M\to\infty\) we have \(K/h\asymp A\log A\), \(\log K\sim\log M\),
\(A\log A/\log M\to c\) and \(\pi(K)h/M=o(M)\).  Hence, \emph{whenever}
\(\theta(K)\le(1+\frac1{8c})K\) and \(M\) is large, \(\log C_M\le(c+\frac12)M\).

\emph{Good scales.}  There are infinitely many such \(M\).  Indeed
\(K(M+1)\le(1+7/M)K(M)\) for large \(M\); given a large \(x\) with
\(\theta(x)\le(1+\frac1{16c})x\) (Lemma~\ref{lem:cheb}), take \(M\) maximal with \(K(M)\le x\);
then \(\theta(K(M))\le\theta(x)\le(1+\frac1{16c})K(M+1)\le(1+\frac1{8c})K(M)\).

\emph{The construction.}  Let \(M\) be a large good scale, \(J=\lfloor e^{(c+2)M}\rfloor\) and
\(F=\{1,\dots,\lfloor e^M\rfloor\}\).  Then \(C_M(|F|+1)\le2e^{(c+3/2)M}<J+1\), and
Lemma~\ref{lem:anchor} gives \(t=j-i\) with \(e^M<t\le e^{(c+2)M}\) and \(\kappa(i)=\kappa(j)\);
thus \(tL\equiv e_p\pmod p\) with \(|e_p|<h\) for every \(p\in\mathcal P\).  Put \(n=tL-h-1\).
Every prime \(p\le M\) divides \(L\mid n+h+1\), and every \(p\in\mathcal P\) divides one of
\(n+2,\dots,n+2h\) by Lemma~\ref{lem:fold} (note \(2h<M<p\)).  Moreover
\[
M<\log t\le\log n\le\log t+\psi(M)\le(c+2+C_1)M=:sM .
\]
Since \(2h<M\le\log n\), every prime \(p\le K\) divides \(\prod_{i\le\log n}(n+i)\), so
\(q(n,\log n)>K\ge AM-1\).

\emph{Conclusion.}  With \(X=\log n\in[M,sM]\),
\[
(1-\varepsilon)\,g(n)\le(1-\varepsilon)\,\frac{sM(\log M+\log s)}{\log\log M}
\le\frac{cM\log M}{\log\log M}-1=AM-1<q(n,\log n)
\]
for large \(M\), because \(c-(1-\varepsilon)s=1+\varepsilon(2+C_1)\ge1\).  As
\(\log n\ge M\to\infty\), this gives infinitely many \(n\). \(\qed\)

\begin{remark}
Two points remove Tao's factor \(\tfrac12\).  The covering block has length
\(2h\asymp\log n/\log\log\log n=o(\log n)\), yet the cost \(\log(p/h)\) of a prime differs from
\(\log(p/\log n)\) only by \(O(\log\log\log\log n)=o(\log\log\log n)\); and the anchored choice
\(t>e^M\) makes \(n\) large at a price \(e^{M}\) that is negligible as \(c\to\infty\).  The same
code was used in \cite{Bae684} for Problem~684.  No prime number theorem is needed: the entropy
bound is only required at the scales actually used, which Lemma~\ref{lem:cheb} supplies.
\end{remark}

\section{Formal verification}

Theorem~\ref{thm}, its \(\limsup\) form, and the Formal Conjectures variant
\texttt{erdos\_457.variants.qnk} are verified in Lean~4 (v4.35.0-rc3) with Mathlib
\cite{Mathlib} (commit \texttt{b84a70d}) at
\url{https://github.com/jbaelaw/erdos457-lean}.  The definition of \(q\) is copied verbatim from
\cite{FC}; the main statement is
\begin{lstlisting}
theorem erdos_457_sharp (ε : ℝ) (hε : 0 < ε) :
    {n : ℕ | (1 - ε) * Real.log n * Real.log (Real.log n) /
        Real.log (Real.log (Real.log n)) ≤ q n (Real.log n)}.Infinite
\end{lstlisting}
and it depends only on the axioms \texttt{propext}, \texttt{Classical.choice} and
\texttt{Quot.sound}.

\section*{Acknowledgements}
I thank the participants of the discussion of Problem~457 on erdosproblems.com, in particular
Terence Tao, whose comment anticipated the constant~\(1\).

\begin{thebibliography}{9}
\small
\bibitem{Bae684}
J.~H.~Bae, \emph{Unbounded logarithmic limsup in Erd\H{o}s Problem 684 via shifted carry
scheduling}, arXiv:2604.23784 (2026).

\bibitem{Bloom457}
T.~F.~Bloom, \emph{Erd\H{o}s problems}, \url{https://www.erdosproblems.com}, Problems 457 and
1181 and their discussion threads.

\bibitem{Erdos1979}
P.~Erd\H{o}s, \emph{Some unconventional problems in number theory}, Acta Math. Acad. Sci.
Hungar. \textbf{33} (1979), 71--80.

\bibitem{FC}
Google DeepMind, \emph{Formal Conjectures},
\url{https://github.com/google-deepmind/formal-conjectures},
\texttt{FormalConjectures/ErdosProblems/457.lean}.

\bibitem{Mathlib}
The mathlib Community, \emph{The Lean mathematical library}, CPP 2020, 367--381.
\end{thebibliography}

\end{document}
